%%%%%%%%%%%%%%%%
%%% gen 18 %%%
%%%%%%%%%%%%%%%%

\Chapter{18}

\Verse{1}
{
	\hJ{וַיֵּרָא אֵלָיו יְהֹוָה בְּאֵלֹנֵי מַמְרֵא וְהוּא יֹשֵׁב פֶּתַח הָאֹהֶל כְּחֹם הַיּוֹם׃}
}
{
	\eJ{
		And YHWH appeared to him at the oaks of Mamre. And he was
		sitting at the tent entrance in the heat of the day,
	}
}
{
	Ok so YHWH appears to him\footnote{Presumably Abraham. YHWH seems mostly focused on Abe these days.}
	near some trees. You know the ones.

	Abe's sitting in the door of his tent. It's hot.
}

\Verse{2}
{
	\hJ{וַיִּשָּׂא עֵינָיו וַיַּרְא וְהִנֵּה שְׁלֹשָׁה אֲנָשִׁים נִצָּבִים עָלָיו וַיַּרְא וַיָּרץ לִקְרָאתָם מִפֶּתַח הָאֹהֶל וַיִּשְׁתַּחוּ אָרְצָה׃}
}
{
	\eJ{
		and he raised his eyes and saw, and here were three people
		standing over him. And he saw and ran toward them from the
		tent entrance and bowed to the ground.
	}
}
{
	Abe looks up and there are three people standing over him.

	He runs over and bows.
%	\footnote{
%		bowed. The word in Hebrew refers to bowing prostrate, nose to the
%		ground.
%	}
}

\Verse{3}
{
	\hJ{וַיֹּאמַר אֲדֹנָי אִם נָא מָצָאתִי חֵן בְּעֵינֶיךָ אַל נָא תַעֲבֹר מֵעַל עַבְדֶּךָ׃}
}
{
	\eJ{
		And he said, “My Lord, if I’ve found favor in your eyes
		don’t pass on from your servant.
	}
}
{

	And Abe says ``My Lord, don't leave me.''
	
	Then uses singular grammar for the three people. Several times.

	\Table[\relax]{llll}{
	\textbf{Hebrew}
		& \textbf{Meaning}
		& \textbf{Actual}
		& \textbf{Expected Plural} \\
	\hline
	\heb{אדני}
		& my lord / my lords / Lord
		& Ambiguous
		& --- \\
	\heb{בעיניך}
		& in your eyes
		& Singular
		& \heb{בעיניכם} \\
	\heb{תעבר}
		& you pass
		& Singular
		& \heb{תעברו} \\
	\heb{עבדך}
		& your servant
		& Singular
		& \heb{עבדכם} \\
	}

%	\footnote{
%		My Lord. Why does Abraham run to the three people and bow to them but
%		then speak to God? Is he meeting God, people, or angels? Several of the
%		biblical stories involving angels contain confusions such as this, i.e.,
%		confusions between when it is the deity and when it is the angel who is
%		speaking or doing something. But it is confusing only so long as we
%		imagine angels as beings who are independent or separate from God. These
%		texts indicate that angels are rather conceived of here as expressions
%		of God’s presence. The consistent biblical conception of God is that God
%		cannot possibly be seen by a human (“A human will not see me and live,”
%		Exod 33:20) and cannot possibly be contained in any known space (“The
%		heavens … will not contain you,” 1 Kings 8:27). God, in this conception,
%		can nonetheless make Himself known to humans by a sort of emanation from
%		the Godhead that is visible to human eyes. It is a hypostasis, a
%		concrete expression of the divine presence, which is otherwise
%		inexpressible to human beings. What the human sees when such a
%		hypostasis is in front of him or her looks like “people,” like a “man.”
%		And the word for such a thing is “angel.” Thus Jacob can encounter an
%		angel and still say “I’ve seen God face-to-face” (Gen 32:31). And
%		Abraham can face three angels and address “my Lord.” And an angel can
%		speak God’s words in first person or can speak about God in third
%		person. All this is so because in some ways an angel is an identifiable
%		thing itself, and in some ways it is merely a representation of divine
%		presence in human affairs. (See The Hidden Face of God, pp. 10–11.)
%	}
}

\Verse{4}
{
	\hJ{יֻקַּח נָא מְעַט מַיִם וְרַחֲצוּ רַגְלֵיכֶם וְהִשָּׁעֲנוּ תַּחַת הָעֵץ׃}
}
{
	\eJ{
		Let a little water be gotten, and wash your feet and relax
		under a tree,
	}
}
{
	Abe then says to YHWH:

	\begin{enumerate}
	\item Let me get you some water.
	\item Wash your feet.%
		\footnote{
		As we learned from J in Genesis 3:8, YHWH has feet.
		Now we learn that YHWH's feet could use a wash.
		}
	\item Rest under a tree.%
		\footnote{
		As we learned from P in Genesis 2:2-3, YHWH sometimes rests.
		Now we learn that YHWH might benefit from resting under a tree, presumably to avoid the sun.
		}
	\end{enumerate}
	
	A strange thing to say to a divine being, but what do I know.

%	\Table{llll}{
%	\textbf{Hebrew} & \textbf{Meaning} & \textbf{Actual} & \textbf{Expected Singular} \\
%	\hline
%	\heb{ורחצו} & wash & Plural & \heb{ורחץ} \\
%	\heb{רגליכם} & your feet & Plural & \heb{רגלך} \\
%	\heb{והשענו} & relax / rest & Plural & \heb{והשען} \\
%	}

}

\Verse{5}
{
	\hJ{וְאֶקְחָה פַת לֶחֶם וְסַעֲדוּ לִבְּכֶם אַחַר תַּעֲבֹרוּ כִּי עַל כֵּן עֲבַרְתֶּם עַל עַבְדְּכֶם וַיֹּאמְרוּ כֵּן תַּעֲשֶׂה כַּאֲשֶׁר דִּבַּרְתָּ׃}
}
{
	\eJ{
		and let me get a bit of bread, and satisfy your heart.
		Afterward you’ll pass on, for that’s why you’ve passed by
		your servant.” And they said, “Do that, as you’ve spoken.”
	}
}
{
	Abe then says to YHWH:

	\begin{enumerate}
	\item Let me get you some bread.
	\item To satisfy your heart.%
		\footnote{
			As we learned from Genesis 6:6, YHWH has a heart, and can be grieved to it.
			Now we learn that YHWH's heart can also be satisfied, and that bread works.
		}
	\item Then you'll leave.
	\item Because that's why you're here.%
		\footnote{YHWH came for bread. Not according to me. According to Abe.}
	\end{enumerate}

	And \textsl{they} say, yeah,

	“Do that. What you said.”
}

\Verse{6}
{
	\hJ{וַיְמַהֵר אַבְרָהָם הָאֹהֱלָה אֶל שָׂרָה וַיֹּאמֶר מַהֲרִי שְׁלֹשׁ סְאִים קֶמַח סֹלֶת לוּשִׁי וַעֲשִׂי עֻגוֹת׃}
}
{
	\eJ{
		And Abraham hurried to the tent, to Sarah, and said,
		“Hurry, three measures of fine flour. Knead it and make
		cakes.”
	}
}
{
	So Abe (now with a different name) runs to his wife (now with a different name) and goes:
	
	``Hurry!''

	He has her make cake.

	I guess that's the bread YHWH came for.
}

\Verse{7}
{
	\hJ{וְאֶל הַבָּקָר רָץ אַבְרָהָם וַיִּקַּח בֶּן בָּקָר רַךְ וָטוֹב וַיִּתֵּן אֶל הַנַּעַר וַיְמַהֵר לַעֲשׂוֹת אֹתוֹ׃}
}
{
	\eJ{
		And Abraham ran to the herd and took a calf, tender and
		good, and he gave it to a servant, and he hurried to
		prepare it.
	}
}
{
	Abe goes and gets a cow. One of the good ones. They have guests.

	He gives it to one of his guys to cook.
}

\Verse{8}
{
	\hJ{וַיִּקַּח חֶמְאָה וְחָלָב וּבֶן הַבָּקָר אֲשֶׁר עָשָׂה וַיִּתֵּן לִפְנֵיהֶם וְהוּא עֹמֵד עֲלֵיהֶם תַּחַת הָעֵץ וַיֹּאכֵלוּ׃}
}
{
	\eJ{
		And he took curds and milk and the calf that he had
		prepared and placed them in front of them, and he was
		standing over them under the tree, and they ate.
	}
}
{
	Abe gives YHWH some cottage cheese and beef.

	And the other two guys too.

	Or the three people.

	Whoever it was.

	The text now notes that Abe is standing \textsl{over} them, \textsl{under} the tree.

	And they ate.
	
	You've gotta love this book.
}

\Verse{9}
{
	\hJ{וַיֹּאמְרוּ אֵלָיו אַיֵּה שָׂרָה אִשְׁתֶּךָ וַיֹּאמֶר הִנֵּה בָאֹהֶל׃}
}
{
	\eJ{
		And they said to him, “Where is Sarah, your wife?” And he
		said, “Here in the tent.”
	}
}
{
	As in Genesis 3:9, and 4:9, and 16:8, YHWH often wants to know where people are.
	
	Abe says ``She's right here.''
}

\Verse{10}
{
	\hJ{וַיֹּאמֶר שׁוֹב אָשׁוּב אֵלֶיךָ כָּעֵת חַיָּה וְהִנֵּה בֵן לְשָׂרָה אִשְׁתֶּךָ וְשָׂרָה שֹׁמַעַת פֶּתַח הָאֹהֶל וְהוּא אַחֲרָיו׃}
}
{
	\eJ{
		And He said, “I shall come back to you at the time of
		life, and, here, Sarah, your wife, will have a son.” And
		Sarah was listening at the tent entrance, which was behind
		him.
	}
}
{
	YHWH says:
	
	I'll be back at the time of life.
	
	And here Sarah, your wife.
	
	She'll have a son.
	
	Sarah's listening.
	
	She's behind YHWH.
}

\Verse{11}
{
	\hJ{וְאַבְרָהָם וְשָׂרָה זְקֵנִים בָּאִים בַּיָּמִים חָדַל לִהְיוֹת לְשָׂרָה אֹרַח כַּנָּשִׁים׃}
}
{
	\eJ{
		And Abraham and Sarah were old, well along in days; Sarah
		had stopped having the way of women.
	}
}
{
	It then reminds us how old they are again.
	
	Sarah has stopped doing that girl thing.\footnote{\aB{.}}
}

\Verse{12}
{
	\hJ{וַתִּצְחַק שָׂרָה בְּקִרְבָּהּ לֵאמֹר אַחֲרֵי בְלֹתִי הָיְתָה לִּי עֶדְנָה וַאדֹנִי זָקֵן׃}
}
{
	\eJ{
		And Sarah laughed inside her and said, “After I’ve become
		worn out am I to have pleasure?! And my lord is old!”
	}
}
{
	Sarah laughs inside her mind.
	
	Because y'know, she's like a hundred.
}

\Verse{13}
{
	\hJ{וַיֹּאמֶר יְהֹוָה אֶל אַבְרָהָם לָמָּה זֶּה צָחֲקָה שָׂרָה לֵאמֹר הַאַף אֻמְנָם אֵלֵד וַאֲנִי זָקַנְתִּי׃}
}
{
	\eJ{
		And YHWH said to Abraham, “Why is this? Sarah laughed,
		saying, ‘Shall I indeed give birth? And I am old!’ Is
		anything too wondrous for YHWH?
	}
}
{
	So YHWH's like ``Why's Sarah laughing?''
	
	``You saying I can't make her pregnant?''
}

\Verse{14}
{
	\hJP{הֲיִפָּלֵא מֵיְהֹוָה דָּבָר לַמּוֹעֵד אָשׁוּב אֵלֶיךָ כָּעֵת חַיָּה וּלְשָׂרָה בֵן׃}
}
{
	\eJP{
		At the appointed time I’ll come back to you, at the time
		of life, and Sarah will have a son.”
	}
}
{
	I'll come back at the appointed time.\footnote{When?}
	
	At life time.\footnote{Thanks, got it.}
	
	And Sarah will have a son.
}

\Verse{15}
{
	\hJ{וַתְּכַחֵשׁ שָׂרָה לֵאמֹר לֹא צָחַקְתִּי כִּי יָרֵאָה וַיֹּאמֶר לֹא כִּי צָחָקְתְּ׃}
}
{
	\eJ{
		And Sarah lied, saying, “I didn’t laugh,” because she was
		afraid. And He said, “No, but you did laugh.”
	}
}
{
	So Sarah's like ``I didn't laugh.''
	
	The text notes that the lies because she's afraid of YHWH.
	
	And YHWH's like ``Yes you did.''
}

\Verse{16}
{
	\hJ{וַיָּקֻמוּ מִשָּׁם הָאֲנָשִׁים וַיַּשְׁקִפוּ עַל פְּנֵי סְדֹם וְאַבְרָהָם הֹלֵךְ עִמָּם לְשַׁלְּחָם׃}
}
{
	\eJ{
		And the people got up from there and gazed at the sight of
		Sodom, and Abraham was going with them to send them off.
	}
}
{
	Anyways the people\footnote{I give up.} get up and go look at Sodom off in the distance.
	
	Abe's like ``Well I don't want to delay you any further,
	you've been so kind for staying so long anyways the exit's
	that way don't forget the leftover cottage cheese and cow.''
}

\Verse{17}
{
	\hJ{וַיהֹוָה אָמָר הַמְכַסֶּה אֲנִי מֵאַבְרָהָם אֲשֶׁר אֲנִי עֹשֶׂה׃}
}
{
	\eJ{
		And YHWH had said, “Shall I conceal what I’m doing from Abraham,
	}
}
{
	And YHWH had\footnote{
		This is that ``noun before verb'' grammar again.
		The one from Genesis 4:1 that means ``had verbed.''
		Like in the past, before this part of the story.
		So this sentence is like a flashback I guess.
	}
	said%
	\aB{\footnote{To who?}}
	
	``Should I lie to Abe about what I'm up to?''
}

\Verse{18}
{
	\hJ{וְאַבְרָהָם הָיוֹ יִהְיֶה לְגוֹי גָּדוֹל וְעָצוּם וְנִבְרְכוּ בוֹ כֹּל גּוֹיֵי הָאָרֶץ׃}
}
{
	\eJ{
		since Abraham will become a big and powerful nation, and
		all the nations of the earth will be blessed through him?
	}
}
{
	``After all,'' continues YHWH talking to some undefined entity,
	
	``he's gonna be big and powerful one day, and bless all the nations.''
}

\Verse{19}
{
	\hJ{כִּי יְדַעְתִּיו לְמַעַן אֲשֶׁר יְצַוֶּה אֶת בָּנָיו וְאֶת בֵּיתוֹ אַחֲרָיו וְשָׁמְרוּ דֶּרֶךְ יְהֹוָה לַעֲשׂוֹת צְדָקָה וּמִשְׁפָּט לְמַעַן הָבִיא יְהֹוָה עַל אַבְרָהָם אֵת אֲשֶׁר דִּבֶּר עָלָיו׃}
}
{
	\eJ{
		For I've chosen him so that he'll command his children and
		his household after him, and they'll keep YHWH's way by
		practicing righteousness and justice, so that YHWH will
		fulfill what He promised Abraham.
	}
}
{

	Ok strap in folks, this one's a doozy.
	
	YHWH says:
	\begin{quote}
	I've chosen him\\
	\hspace*{2em}so that\\
	\hspace*{4em}he'll command his children\\
	\hspace*{4em}and his household after him,\\
	\hspace*{4em}and they'll keep YHWH's way\\
	\hspace*{6em}by practicing righteousness and justice,\\
	\hspace*{2em}so that\\
	\hspace*{4em}YHWH will fulfill what He promised Abraham.
	\end{quote}

	Ok so as far as I can tell, the two ``so that'' clauses are parallel, not nested.
	
	So it's not

	\begin{quote}
	I've chosen him\\
	\hspace*{2em}so that\\
	\hspace*{4em}he'll do X\\
	\hspace*{4em}so that\\
	\hspace*{6em}I (YHWH) will do Y.
	\end{quote}

	Rather, it's

	\begin{quote}
	I've chosen him\\
	\hspace*{2em}so that\\
	\hspace*{4em}he'll do X\\
	\hspace*{2em}so that\\
	\hspace*{4em}I (YHWH) will do Y.
	\end{quote}

	\aB{Why?}

	I'm not gonna explain. But there's reasons. Let's keep moving.

	\aC{
		Biblical Hebrew commonly repeats \heb{לְמַעַן} to introduce
		multiple parallel purpose clauses. Here, the second \heb{לְמַעַן}
		appears to return to the main verb \heb{יְדַעְתִּיו} rather than
		depending on the first \heb{לְמַעַן}, making the two purposes
		grammatically parallel.
	}

	\aB{Thanks.}

	No problem.

	Ok so...
	
	``I've chosen him'' here presumably means ``chosen for the covenant.''
	
	That is, ``I've chosen him so that'' means ``I've made a promise to him so that''...

	In which case, it's this:

	\begin{quote}
	I've made a promise to A\\
	\hspace*{2em}so that\\
	\hspace*{4em}A will do stuff I want\\
	\hspace*{2em}(and also)\\
	\hspace*{2em}so that\\
	\hspace*{4em}I (YHWH) will keep my promise to A.
	\end{quote}

	So basically the point of all that is that YHWH chose Abe.

	So that:\\
	(1) Abe would do stuff, and\\
	(2) YHWH would keep his promise to Abe.

	\aB{What?}

	Let's just focus on item (2) in case that seemed clear.

	Ok so...

	\begin{quote}
	YHWH made a promise to Abe.

	So that YHWH would keep his promise to Abe.%
	\footnote{And the other reason from (2), omitted this time for, um, clarity.}
	\end{quote}

	I think that's what he just said.

	But don't quote me on that.	
}

\Verse{20}
{
	\hJ{וַיֹּאמֶר יְהֹוָה זַעֲקַת סְדֹם וַעֲמֹרָה כִּי רָבָּה וְחַטָּאתָם כִּי כָבְדָה מְאֹד׃}
}
{
	\eJ{
		And YHWH said, “The cry of Sodom and Gomorrah: how great
		it is. And their sin: how very heavy it is.
	}
}
{
	Ok so YHWH mentions that Sodom and Gomorrah are:
	
	\begin{enumerate}
	\item Sinning.
	\item Crying.
	\end{enumerate}

	It doesn't say why they're crying.
	
	We'll see why in a bit.
}

\Verse{21}
{
	\hJ{אֵרְדָה נָּא וְאֶרְאֶה הַכְּצַעֲקָתָהּ הַבָּאָה אֵלַי עָשׂוּ כָּלָה וְאִם לֹא אֵדָעָה׃}
}
{
	\eJ{
		Let me go down, and I’ll see if they’ve done, all told,
		like the cry that has come to me. And if not, let me
		know.”
	}
}
{
	I think what he's saying here is:
	
	``I'll go down and see if it's really as bad as I've heard. If not, I'll find out.''
	
	The ``let me know'' in the translation is more like ``let it be the case that I know.''

	\aC{
		The Hebrew \heb{אֵדָעָה} is a first-person cohortative, a verbal form
		often expressing resolve, intention, or desire. Here it means
		something like ``I'll know,'' ``I'll find out,'' or ``let me find out,''
		not the modern English sense of the idiom ``let me know.'' 
	}
	
	Point being: I don't think he's asking himself to tell himself what happened.
	
	Then again, there's sort of three YHWHs in this scene.
	
	So don't quote me on that.

}

\Verse{22}
{
	\hJ{וַיִּפְנוּ מִשָּׁם הָאֲנָשִׁים וַיֵּלְכוּ סְדֹמָה וְאַבְרָהָם עוֹדֶנּוּ עֹמֵד לִפְנֵי יְהֹוָה׃}
}
{
	\eJ{
		And the people turned from there and went to Sodom. And
		Abraham was still standing before YHWH,
	}
}
{
	Damn, just when I thought we knew what was going on here.

	Anyways the ``three people'' from 18:2 are now ``two people and YHWH.''	

	Many\footnote{``Many'' as in ``three,'' in the sense of the current chapter.}
	scholars explain this by saying that angels are a ``hypostasis.''%
	\footnote{
		Like watching TV. You saw the football game, but also not.
		It makes sense but I'll spare you. The football game is God.
	}

	Nevertheless, I think it's pretty indisputable that the author\footnote{J, the \textsc{Goat}.} is up to something.
}

\Verse{23}
{
	\hJ{וַיִּגַּשׁ אַבְרָהָם וַיֹּאמַר הַאַף תִּסְפֶּה צַדִּיק עִם רָשָׁע׃}
}
{
	\eJ{
		and Abraham came over and said, “Will you also annihilate
		the virtuous with the wicked?
	}
}
{
	Though it hasn't exactly been established yet that ``annihilation'' is YHWH's intention here,

	Abe has known YHWH for like 24 years at this point.

	So I guess he knows him well enough to infer.
}

\Verse{24}
{
	\hJ{אוּלַי יֵשׁ חֲמִשִּׁים צַדִּיקִם בְּתוֹךְ הָעִיר הַאַף תִּסְפֶּה וְלֹא תִשָּׂא לַמָּקוֹם לְמַעַן חֲמִשִּׁים הַצַּדִּיקִם אֲשֶׁר בְּקִרְבָּהּ׃}
}
{
	\eJ{
		Maybe there are fifty virtuous people within the city.
		Will you also annihilate and not sustain the place for the
		fifty virtuous who are in it?
	}
}
{
	Abe's like ``What if there's fifty good people?''
}

\Verse{25}
{
	\hJ{חָלִלָה לְּךָ מֵעֲשֹׂת כַּדָּבָר הַזֶּה לְהָמִית צַדִּיק עִם רָשָׁע וְהָיָה כַצַּדִּיק כָּרָשָׁע חָלִלָה לָּךְ הֲשֹׁפֵט כּל הָאָרֶץ לֹא יַעֲשֶׂה מִשְׁפָּט׃}
}
{
	\eJ{
		Far be it from you to do a thing like this, to kill
		virtuous with wicked—and it will be the same for the
		virtuous and the wicked—far be it from you. Will the judge
		of all the earth not do justice?”
	}
}
{
	``You wouldn't kill \textsc{fifty} people just for being nearby, would you?''

%	\footnote{
%		Far be it from you to do a thing like this. This is the first time in
%		the Bible that a human questions a divine decision. Moses will take this
%		even farther on at least three occasions. Concerning the exchange
%		between God and Abraham, see the comment on Gen 22:3.
%	}
}

\Verse{26}
{
	\hJ{וַיֹּאמֶר יְהֹוָה אִם אֶמְצָא בִסְדֹם חֲמִשִּׁים צַדִּיקִם בְּתוֹךְ הָעִיר וְנָשָׂאתִי לְכל הַמָּקוֹם בַּעֲבוּרָם׃}
}
{
	\eJ{
		And YHWH said, “If I find in Sodom fifty virtuous people
		within the city, then I’ll sustain the whole place for
		their sake.”
	}
}
{
	And YHWH does some quick math in his head and goes ``No. I wouldn't do that.''
}

\Verse{27}
{
	\hJ{וַיַּעַן אַבְרָהָם וַיֹּאמַר הִנֵּה נָא הוֹאַלְתִּי לְדַבֵּר אֶל אֲדֹנָי וְאָנֹכִי עָפָר וָאֵפֶר׃}
}
{
	\eJ{
		And Abraham answered, and he said, “Here I’ve undertaken
		to speak to my Lord, and I’m dust and ashes.
	}
}
{
	Abe says a thing about being dust and ashes.
	
	On with the scene.
}

\Verse{28}
{
	\hJ{אוּלַי יַחְסְרוּן חֲמִשִּׁים הַצַּדִּיקִם חֲמִשָּׁה הֲתַשְׁחִית בַּחֲמִשָּׁה אֶת כּל הָעִיר וַיֹּאמֶר לֹא אַשְׁחִית אִם אֶמְצָא שָׁם אַרְבָּעִים וַחֲמִשָּׁה׃}
}
{
	\eJ{
		Maybe the fifty virtuous people will be short by five.
		Will you destroy the whole city for the five?” And He
		said, “I won’t destroy if I find forty-five there.”
	}
}
{
	Abe doesn't say ``forty five'' here.
	
	He does a ``dollar ninety-nine'' sort of thing.\aC{\footnote{In reverse.}}
	
	YHWH says ``No, I won't kill everybody if there's forty-five good ones.''
	
	Note: YHWH's answer dispenses with Abe's marketing speak.
}

\Verse{29}
{
	\hJ{וַיֹּסֶף עוֹד לְדַבֵּר אֵלָיו וַיֹּאמַר אוּלַי יִמָּצְאוּן שָׁם אַרְבָּעִים וַיֹּאמֶר לֹא אֶעֱשֶׂה בַּעֲבוּר הָאַרְבָּעִים׃}
}
{
	\eJ{
		And he went on again to speak to Him and said, “Maybe
		forty will be found there.” And He said, “I won’t do it
		for the sake of the forty.”
	}
}
{
	Abe's like ``Ok forty.''

	YHWH's like ``No Abe, I won't kill everybody if there's forty good ones.''

}

\Verse{30}
{
	\hJ{וַיֹּאמֶר אַל נָא יִחַר לַאדֹנָי וַאֲדַבֵּרָה אוּלַי יִמָּצְאוּן שָׁם שְׁלֹשִׁים וַיֹּאמֶר לֹא אֶעֱשֶׂה אִם אֶמְצָא שָׁם שְׁלֹשִׁים׃}
}
{
	\eJ{
		And he said, “May my Lord not be angry, and let me speak.
		Maybe thirty will be found there.” And He said, “I won’t
		do it if I find thirty there.”
	}
}
{
	So Abe's like ``Don't get mad. Let me talk.''

	``...Thirty?''

	YHWH's like ``No Abe, I won't kill everybody if there's thirty.''
}

\Verse{31}
{
	\hJ{וַיֹּאמֶר הִנֵּה נָא הוֹאַלְתִּי לְדַבֵּר אֶל אֲדֹנָי אוּלַי יִמָּצְאוּן שָׁם עֶשְׂרִים וַיֹּאמֶר לֹא אַשְׁחִית בַּעֲבוּר הָעֶשְׂרִים׃}
}
{
	\eJ{
		And he said, “Here I’ve undertaken to speak to my Lord.
		Maybe twenty will be found there.” And He said, “I won’t
		destroy for the sake of the twenty.”
	}
}
{
	So Abe's like ``Look, please, I got up the courage to talk to you.''
	
	``Twenty?''

	YHWH's like ``No Abe, not for twenty.''
	
	Side Note: Haven't mentioned this yet, but Abe's nephew lot is down there in Sodom.
	
	Abe doesn't mention that.
	
	He's just sort of asking abstractly.
}

\Verse{32}
{
	\hJ{וַיֹּאמֶר אַל נָא יִחַר לַאדֹנָי וַאֲדַבְּרָה אַךְ הַפַּעַם אוּלַי יִמָּצְאוּן שָׁם עֲשָׂרָה וַיֹּאמֶר לֹא אַשְׁחִית בַּעֲבוּר הָעֲשָׂרָה׃}
}
{
	\eJ{
		And he said, “May my Lord not be angry, and let me speak
		just this one time. Maybe ten will be found there.” And He
		said there, “I won’t destroy for the sake of the ten.”
	}
}
{
	So Abe's like ``Please, don't get mad, just let me say one more thing...''
	
	``Ten?''
	
	``Maybe there's ten.''
	
	And YHWH's like ``No Abe, not for ten.''

}

\Verse{33}
{
	\hJ{וַיֵּלֶךְ יְהֹוָה כַּאֲשֶׁר כִּלָּה לְדַבֵּר אֶל אַבְרָהָם וְאַבְרָהָם שָׁב לִמְקֹמוֹ׃}
}
{
	\eJ{
		And YHWH went when He had finished speaking to Abraham,
		and Abraham went back to his place.
	}
}
{
	YHWH leaves, taking with him the number he silently computed in his head back in verse 26, smiling toward Sodom.
	
	Abe goes home.
	
	End of chapter.
	
	Up next, the destruction of Sodom.
}
